\subsection{同型模的描述}

与前一小节一样，A 表示一个环，S 表示一个单左 A-模，D 表示体 \(\operatorname{End}_{\mathrm{A}}(\mathrm{S})^{\circ}\)。我们把 S 看作一个 (A,D)-双模。

定义 3. —— 设 M 是一个 S 型同型 A-模。M 关于 S 的一个描述是一个偶 \((V, \alpha)\)，其中 V 是体 D 上的一个左向量空间，而 \(\alpha : S \otimes_{D} V \to M\) 是一个 A-模同构。

每个 S 型同型 A-模 M 都有一个典范描述：它就是偶 \((\mathrm{Hom}_{\mathrm{A}}(\mathrm{S},\mathrm{M}),\alpha_{\mathrm{M}})\)，其中 \(\alpha_{M}:S\otimes_{D}\mathrm{Hom}_{\mathrm{A}}(\mathrm{S},\mathrm{M})\to\mathrm{M}\) 是由 \(\alpha_{\mathrm{M}}(s\otimes f)=f(s)\)（VIII，第 62 页，命题 3, a)）刻画的 A-模同构。

定理 2. —— 设 M 是一个 S 型同型 A-模，\((V,\alpha)\) 是 M 的一个描述。用 \(\mathcal{D}_{\mathrm{D}}(\mathrm{V})\) 表示 V 的 D-线性子空间组成的集合（按包含关系排序），用 \(\mathcal{D}_{\mathrm{A}}(\mathrm{M})\) 表示 M 的 A-子模组成的集合。对每个 \(W \in \mathcal{D}_{\mathrm{D}}(\mathrm{V})\)，把 A-模 \(S \otimes_{D} W\) 与它在 \(S \otimes_{D} V\) 中的典范像等同起来。

\begin{enumerate}
\def\labelenumi{\alph{enumi})}
\item
  映射 \(W \mapsto \alpha(S \otimes_{D} W)\) 是从 \(\mathcal{D}_{\mathrm{D}}(\mathrm{V})\) 到 \(\mathcal{D}_{\mathrm{A}}(\mathrm{M})\) 的一个有序集同构。
\item
  逆同构把 M 的一个子模 N 映到 V 的由这样的元素 v 组成的线性子空间：\(\alpha(s \otimes v)\)（对每个 \(s \in S\)）属于 N。
\end{enumerate}

对 \(\mathrm{W} \in \mathcal{D}_{\mathrm{D}}(\mathrm{V})\)，令 \(\varphi(\mathrm{W}) = \alpha(\mathrm{S} \otimes_{\mathrm{D}} \mathrm{W})\)。对 \(\mathrm{N} \in \mathcal{D}_{\mathrm{A}}(\mathrm{M})\)，用 \(\psi(\mathrm{N})\) 表示这样的元素 \(v \in \mathrm{V}\) 组成的集合：\(\alpha(s \otimes v) \in \mathrm{N}\) 对每个 \(s \in \mathrm{S}\) 成立。这就定义了两个映射 \(\varphi: \mathcal{D}_{\mathrm{D}}(\mathrm{V}) \to \mathcal{D}_{\mathrm{A}}(\mathrm{M})\) 与 \(\psi: \mathcal{D}_{\mathrm{A}}(\mathrm{M}) \to \mathcal{D}_{\mathrm{D}}(\mathrm{V})\)。它们显然都是递增的。

设 N 是 M 的一个子模。它是 S 型同型的（VIII，第 61 页，命题 2）。令 \(W = \psi(N)\)。由 VIII，第 62 页的命题 3, b)，A-线性映射 \(h : S \to M\) 就是映射 \(s \mapsto \alpha(s \otimes v)\)，其中 v 取遍 V。那些像包含于 N 的是映射 \(s \mapsto \alpha(s \otimes w)\)，其中 w 取遍 W；由于 N 是 S 型同型的，它们的像生成 N。于是我们有 \(\alpha(S \otimes_{D} W) = N\)，即 \(\varphi(\psi(N)) = N\)。这证明 \(\varphi \circ \psi\) 是 \(\mathcal{D}_{\mathrm{A}}(\mathrm{M})\) 上的恒等映射。特别地，\(\varphi\) 是满射而 \(\psi\) 是单射。

为完成证明，只需证映射 \(\varphi\) 是单射。设 W 与 \(W'\) 是 V 的线性子空间，使得 \(\varphi(\mathrm{W}) = \varphi(\mathrm{W}')\)。向量空间 \(S \otimes_{D} W\) 与 \(S \otimes_{D} W'\) 作为 \(S \otimes_{D} V\) 的线性子空间是一致的。取 D-向量空间 S 上的一个非零线性型 f，并设 \(g : S \otimes_{D} V \to V\) 是由 \(g(s \otimes v) = f(s)v\) 定义的群同态。我们有 \(W = g(S \otimes_{D} W) = g(S \otimes_{D} W') = W'\)，故 \(\varphi\) 是单射。

注 1. —— 设 M 是一个 S 型同型 A-模，\((V, \alpha)\) 是 M 的一个描述。则 M 长度有限当且仅当 V 是有限维的，且在这一情形，我们有关系式

\[
\mathrm{long} _ {\mathrm{A}} (\mathrm{M}) = \mathrm{dim} _ {\mathrm{D}} (\mathrm{V}).
\]

推论 1. —— 设 M 是一个 S 型同型 A-模。对 M 的每个 A-子模 N，把 \(\operatorname{Hom}_{\mathrm{A}}(\mathrm{S},\mathrm{N})\) 与 \(\operatorname{Hom}_{\mathrm{A}}(\mathrm{S},\mathrm{M})\) 的由像包含于 N 的映射组成的线性子空间等同起来。

\begin{enumerate}
\def\labelenumi{\alph{enumi})}
\item
  映射 \(N \mapsto \operatorname{Hom}_{\mathrm{A}}(\mathrm{S}, \mathrm{N})\) 是从 \(\mathcal{D}_{\mathrm{A}}(\mathrm{M})\) 到 \(\mathcal{D}_{\mathrm{D}}(\operatorname{Hom}_{\mathrm{A}}(\mathrm{S}, \mathrm{M}))\) 的一个有序集同构。
\item
  逆双射把线性子空间 \(\mathrm{W}\)（\(\operatorname{Hom}_{\mathrm{A}}(\mathrm{S}, \mathrm{M})\) 的）映到子模 \(\sum_{h \in \mathrm{W}} h(\mathrm{S})\)（\(\mathrm{M}\) 的）。
\end{enumerate}

这是当我们取 \((\mathrm{V},\alpha)\) 为 M 的典范描述时定理 2 的一个重新表述。

推论 2. —— 设 V 是 D 上的一个左向量空间，F 是 V 的自同态组成的一个集合。\(S \otimes_{D} V\) 的一个 A-子模在所有自同态

\(1_{\mathrm{S}} \otimes u\)（其中 \(u\) 取遍 \(\mathcal{F}\)）下稳定，当且仅当它形如 \(\mathrm{S} \otimes_{\mathrm{D}} \mathrm{W}\)，其中 \(\mathrm{W}\) 是 \(\mathrm{V}\) 的一个在所有属于 \(\mathcal{F}\) 的自同态下稳定的线性子空间。

事实上，由定理 2，\(S \otimes_{D} V\) 的每个 A-子模 N 都等于 \(S \otimes_{D} W\)，其中 W 是 V 的由这样的元素 v 组成的线性子空间：\(s \otimes v\)（对每个 \(s \in S\)）属于 N。

定理 3. —— 设 M 与 M' 是 S 型同型 A-模。设 \((V, \alpha)\) 与 \((V', \alpha')\) 分别是 M 与 M' 的描述。对任一 D-线性映射 \(f : V \to V'\)，用 \(\widetilde{f} : M \to M'\) 表示使下图交换的唯一的 A-线性映射：

\[
\begin{array}{c} \mathrm{S} \otimes_ {\mathrm{D}} \mathrm{V} \xrightarrow {\alpha} \mathrm{M} \\ \Biggl \downarrow_ {1 _ {\mathrm{S}} \otimes f} \\ \mathrm{S} \otimes_ {\mathrm{D}} \mathrm{V} ^ {\prime} \xrightarrow {\alpha^ {\prime}} \mathrm{M} ^ {\prime}. \end{array}\tag{9}
\]

映射 \(f \mapsto \widetilde{f}\)（从 \(\mathrm{Hom}_{\mathrm{D}}(\mathrm{V},\mathrm{V}')\) 到 \(\mathrm{Hom}_{\mathrm{A}}(\mathrm{M},\mathrm{M}')\)）是一个群同构。

只需证明 Z-线性映射 \(u \mapsto 1_{S} \otimes u\)（从 \(\mathrm{Hom}_{\mathrm{D}}(\mathrm{V}, \mathrm{V}')\) 到 \(\mathrm{Hom}_{\mathrm{A}}(\mathrm{S} \otimes_{\mathrm{D}} \mathrm{V}, \mathrm{S} \otimes_{\mathrm{D}} \mathrm{V}')\)）是双射。用第 3 目的记号（应用于 (A,D)-双模 S），这相当于证明映射

\[
\mathcal {T}: \mathrm{Hom} _ {\mathrm{D}} (\mathrm{V}, \mathrm{V} ^ {\prime}) \longrightarrow \mathrm{Hom} _ {\mathrm{A}} (\mathcal {T} (\mathrm{V}), \mathcal {T} (\mathrm{V} ^ {\prime}))
\]

是双射。但由 VIII，第 60 页的注 1，由于伴随同构（VIII，第 59 页）是双射，这相当于证明把 u 映为 \(\beta_{V'} \circ u\) 的映射是双射，而这由 D-线性映射 \(\beta_{V'}\) 是双射这一事实得出（VIII，第 62 页，命题 3, b)）。

我们保留定理 3 的记号。设 \(\mathrm{M}''\) 是一个 S 型同型 A-模，\((\mathrm{V}'',\alpha '')\) 是 \(\mathrm{M}''\) 的一个描述。对每个 \(f\in \mathrm{Hom}_{\mathrm{D}}(\mathrm{V},\mathrm{V}')\) 与每个 \(g\in \mathrm{Hom}_{\mathrm{D}}(\mathrm{V}',\mathrm{V}'')\)，我们有 \(\widetilde{g\circ f} = \widetilde{g}\circ \widetilde{f}\)。特别地，对 \(\mathrm{M} = \mathrm{M}'\)、\(\mathrm{V} = \mathrm{V}'\)、\(\alpha = \alpha '\)，映射 \(f\mapsto \widetilde{f}\)（从 \(\mathrm{End}_{\mathrm{D}}(\mathrm{V})\) 到 \(\mathrm{End}_{\mathrm{A}}(\mathrm{M})\)）是一个环同构。

注 2. —— 设 M 是一个 S 型同型 A-模，\((V,\alpha)\) 是 M 的一个描述。设 B 是环 \(\mathrm{End}_{\mathrm{A}}(\mathrm{M})^{\circ}\) 的一个子环。从 \(\mathrm{End}_{\mathrm{D}}(\mathrm{V})^{\circ}\) 到 \(\mathrm{End}_{\mathrm{A}}(\mathrm{M})^{\circ}\) 的环同构赋给 V 一个 \((D,B)\) -双模结构，使 \(\alpha\) 是 \((A,B)\) -双模的一个同构。存在一个从 V 的 \((D,B)\) -双子模组成的集合（按包含关系排序）到 M 的 (A,B)-双子模组成的集合的同构（VIII，第 62 页，定理 2 及 VIII，第 63 页，推论 2）。

推论. —— 设 M 与 M' 是 S 型同型 A-模。映射 \(u \mapsto \operatorname{Hom}(1_{\mathrm{S}}, u)\)（从 \(\operatorname{Hom}_{\mathrm{A}}(\mathrm{M}, \mathrm{M}')\) 到 \(\operatorname{Hom}_{\mathrm{D}}(\operatorname{Hom}_{\mathrm{A}}(\mathrm{S}, \mathrm{M}), \operatorname{Hom}_{\mathrm{A}}(\mathrm{S}, \mathrm{M}'))\)）是一个群同构。当 M 等于 \(M'\) 时，它是从 \(\operatorname{End}_{\mathrm{A}}(\mathrm{M})\) 到 \(\operatorname{End}_{\mathrm{D}}(\operatorname{Hom}_{\mathrm{A}}(\mathrm{S}, \mathrm{M}))\) 的一个环同构。

由于 VIII，第 59 页的图 (I) 的交换性，这个推论由定理 3（应用于 M 与 M' 的典范描述）得出。
% semantic-fragments: legacy-input-seam
\relax{}
